% checklist_fail.tex — 故意埋错的迷你论文 fixture（Wave6-A）
\documentclass{article}
\usepackage{booktabs}
\begin{document}

\section{摘要}
针对问题1，本文采用优化方法求解，得到结果。关键词：优化

\section{引言}
\textbf{问题1}：求解最小成本。

\section{模型假设}
假设1：稳态。
假设2：正态扰动。
假设3：线性关系。

\section{符号说明}
\begin{table}[htbp]
  \centering
  \caption{符号}
  \label{tab:sym}
  \begin{tabular}{cl}
    \toprule
    符号 & 含义 \\
    \midrule
    \bottomrule
  \end{tabular}
\end{table}

\section{模型求解}
经过求解得到结果。
\begin{figure}[htbp]
  \centering
  \includegraphics[width=0.99\textwidth]{big.pdf}
  \caption{A Big Figure}
  \label{fig:big}
\end{figure}
如图\ref{fig:big}所示，\textbf{本步}采用贪心策略。
\begin{table}[htbp]
  \centering
  \caption{未被引用的表}
  \label{tab:orphan}
  \begin{tabular}{lc}
    \toprule
    A & 1 \\
    \bottomrule
  \end{tabular}
\end{table}

\section{模型的评价}
优点1：求解快。
缺点1：未考虑扰动。
缺点2：精度不足。

\begin{thebibliography}{9}
\bibitem{a} 张三. Some paper[J]. J, 2023. DOI: 10.1234/abc.5678.
\end{thebibliography}

\end{document}
